% !TEX program = xelatex

% Le PDF produit doit être convertit en PNG. Pour l'instant, je le fais avec Gimp.

\documentclass[margin=1mm]{standalone}

\usepackage{fontspec}
\usepackage{xunicode}
\usepackage{xcolor}
\usepackage{tikz}
\usetikzlibrary{arrows.meta, positioning}
\setmainfont[Ligatures=TeX]{Source Sans Pro}


\begin{document}

\begin{tikzpicture}[
  % Arrow style
  arrow/.style={->,>=stealth},
  % Boxes style
  box/.style={ rectangle, draw, minimum height=10mm, align=center, inner sep=6pt, text width=8mm, fill=white},
  txt/.style={ align=center, inner sep=6pt},
]

  % Blocks
  \node[txt]         (x) at (  0,  1) {$x$};
  \node[box]         (h) at (  2,  1) {$h$};
  \node[draw,circle] (s) at (  4,  1) {$+$};
  \node[txt]         (b) at (  4,2.5) {$b$};
  \node[txt]         (y) at (4.8,1.2) {$y$};
  \node[box]         (g) at (  6,  1) {$g$};
  \node[txt]         (e) at (  8,  1) {$\hat{x}$};

  % Captions
  \node at (2,0) {degradation};
  \node at (6,0) {restoration};

  % Arrows
  \draw [arrow] (x)--(h);
  \draw [arrow] (h)--(s);
  \draw [arrow] (b)--(s);
  \draw [arrow] (s)--(g);
  \draw [arrow] (g)--(e);

\end{tikzpicture}

\end{document}